% SPDX-License-Identifier: LPPL-1.3c
% Copyright (C) 2026 Christophe Jorssen
% Author and Current Maintainer: Christophe Jorssen <christophe.jorssen@gmail.com>
% LPPL maintenance status: maintained
% This file is part of luafemm. See LICENSE and LICENSES.md for its terms.
% Task-led tutorial with executed setup and source/result blocks.
\clearpage
\section{Tutorial: Kamil looks for a defect without cutting the tube}
\label{sec:tutorial-ndt}

Kamil has a copper tube, a small coil and a new question for his students.
Can the coil tell him that some metal is missing from the tube wall? He would
like to show the experiment in three stages: a drawing of the apparatus,
a field picture explaining the interaction, and the signal that an instrument
could actually measure. A coloured defect in a drawing is easy to find.
The more interesting part is finding it in the coil's electrical response.

We shall build an encircling-coil inspection of a tube with a circumferential
groove. We then move the probe, compare each position with a sound reference,
and plot the two components of the impedance change. The three worked cases
below declare their own styles and shared setup. They use ordinary
TikZ paths, the existing |femm| keys and PGFPlots; no document-side Lua,
native solver or external data file is required.

\subsection{First decide what the defect means}
A meridional $(r,z)$ section of a circular tube is a rectangle. Revolving a
small rectangular recess in its outer edge makes an \emph{annular groove}.
This is a useful model of circumferential material loss. It is also compatible
with the axisymmetric harmonic solver: the coil, specimen and defect all
retain rotational symmetry throughout the scan.

A notch on just one side of the tube, an off-centre probe or a longitudinal
slit would break that symmetry. Drawing one of those in this meridional
picture would still revolve it into a ring. Kamil therefore labels his defect
as a groove and keeps the coil coaxial. The remaining azimuthal currents still
form closed circles; the groove changes the available conducting volume and
its coupling to the coil. It does not force these circles to detour around a
local three-dimensional crack.

The apparatus uses the following dimensions. Coordinates and local mesh
spacings are in millimetres, while conductivity is in SI units.
\begin{center}
\begin{tabular}{ll}
\toprule Setting & Value\\
\midrule
Tube inner / outer radius & $5/8\,\mathrm{mm}$\\
Tube axial extent & $-25\leq z\leq25\,\mathrm{mm}$\\
Groove width / centre & $3\,\mathrm{mm}$ / $z=0$\\
Groove depths to compare & $0.5$ and $1.5\,\mathrm{mm}$\\
Coil radial extent & $9\leq r\leq11\,\mathrm{mm}$\\
Coil axial extent & $s-1.5\leq z\leq s+1.5\,\mathrm{mm}$\\
Coil turns / copper fill & $60$ / $0.6$\\
Excitation & $f=2\,\mathrm{kHz}$, $\widehat I=1\,\mathrm A$ peak\\
Copper & $\mu_r=1$, $\sigma=58\,\mathrm{MS/m}$\\
Air-domain truncation & $0\leq r\leq45$, $-45\leq z\leq45\,\mathrm{mm}$\\
\bottomrule
\end{tabular}
\end{center}
In the apparatus drawing, two orange rectangles are opposite sections of the
\emph{same} encircling winding. The solver's meridional drawing needs only the
positive-radius rectangle. The transverse view shows that the missing metal
extends around the whole circumference. This first drawing is purely
geometrical; the following two examples calculate the fields and signals.
\luafemmexamplebegin{tutorial-ndt-setup}

% manual example: tutorial-ndt-setup/1

\begin{codeexample}[setup code]
\usetikzlibrary{patterns,arrows.meta}
\definecolor{ndtblue}{RGB}{25,83,112}
\definecolor{ndtred}{RGB}{192,69,44}
\end{codeexample}

Apply these styles to the physical paths before requesting the field.

% manual example: tutorial-ndt-setup/2

\begin{codeexample}[width=0pt]
\begin{tikzpicture}[x=1.6mm,y=1.6mm,font=\small,
    dimension/.style={gray!80,{Stealth[length=1.4mm]}-{Stealth[length=1.4mm]}},
    leader/.style={line width=.5pt,-{Stealth[length=1.5mm]}},
    wall/.style={draw=ndtblue,fill=ndtblue!12,line width=.6pt},
    coil/.style={draw=orange!80!black,fill=orange!25,line width=.7pt}]
  \node[anchor=west,font=\bfseries,text=ndtblue] at (-25,21)
    {A moving probe, a fixed circumferential defect};

  \filldraw[wall] (-25,5) rectangle (25,8);
  \filldraw[wall] (-25,-8) rectangle (25,-5);
  \foreach \Side in {-1,1}{
    \fill[pattern=north east lines,pattern color=ndtblue!25]
      (-25,{\Side*5}) rectangle (25,{\Side*8});
    \filldraw[fill=white,draw=ndtred,line width=.8pt]
      (-1.5,{\Side*6.5}) rectangle (1.5,{\Side*8});
    \filldraw[coil] (-9.5,{\Side*9}) rectangle (-6.5,{\Side*11});
    \draw[orange!70,densely dashed]
      (6.5,{\Side*9}) rectangle (9.5,{\Side*11});
  }
  \draw[gray,dash dot,-Stealth] (-27,0)--(28,0) node[right] {$z$};
  \draw[gray,densely dashed] (0,-10)--(0,11);
  \draw[orange!80!black,thick,-Stealth] (-12,15)--(12,15)
    node[midway,above] {Scan position $s$};
  \draw[leader,orange!80!black] (-21,12) node[above] {60 turns}--(-8,10);
  \draw[leader,ndtred] (6,-18) node[below] {Groove, width 3 mm}--(.5,-7);
  \draw[dimension] (-25,-13)--(25,-13)
    node[midway,fill=white,inner sep=2pt] {50 mm};
  \node[text=ndtblue,fill=white,inner sep=1pt] at (16,6.5) {Copper};
  \node[fill=white,inner sep=1pt] at (15,2) {Air bore};

  \begin{scope}[shift={(44,0)}]
    \fill[ndtred!8] (0,0) circle[radius=8];
    \draw[ndtred,densely dashed] (0,0) circle[radius=8];
    \filldraw[wall] (0,0) circle[radius=6.5];
    \filldraw[fill=white,draw=ndtblue] (0,0) circle[radius=5];
    \draw[dimension] (0,0)--(0,5)
      node[midway,right,fill=white,inner sep=1pt] {5};
    \draw[leader,ndtred] (7,12) node[above] {Removed metal}--(35:7.3);
    \draw[dimension] (-8,-11)--(8,-11)
      node[midway,below] {Sound diameter: 16 mm};
    \node[align=center,text=ndtblue] at (0,-18)
      {Transverse section at $z=0$\\Groove depth: 1.5 mm};
  \end{scope}
  \node[anchor=west,font=\footnotesize] at (-25,-22)
    {Copper: $\sigma=58$ MS/m, $\mu_r=1$
     Excitation: $f=2$ kHz, $\widehat I=1$ A (peak)};
\end{tikzpicture}
\end{codeexample}

\luafemmexampleend{tutorial-ndt-setup}

\subsection{Prepare a model that is inexpensive to edit}
Kamil starts with a coarse field mesh while placing labels and deciding how
much of the apparatus to show. The far air needs fewer triangles than the
conducting wall or the narrow groove. At $2\,\mathrm{kHz}$, copper's standard
skin depth is approximately
\[
 \delta=\sqrt{\frac{2}{\omega\mu_0\sigma}}\simeq1.48\,\mathrm{mm}.
\]
He can try |\Mesh=.6| while composing the example, then use |.4| for the
published pictures and compare with |.25| before interpreting small signal
changes. The groove receives half this spacing. These are resolution choices,
not a built-in adaptive skin-depth estimator.

He puts the common physical settings in a TikZ style:
\begin{codeexample}[code only]
\def\Mesh{.4}\def\Frequency{2000}
\pgfmathsetmacro{\NearMesh}{1.5*\Mesh}
\pgfmathsetmacro{\GrooveMesh}{\Mesh/2}
\tikzset{ndt/model/.style={
  femm/problem={type=axisymmetric,frequency=\Frequency,
    xmax=45,ymin=-45,ymax=45,mesh size=2},
  femm/circuit={name=probe,current=1},
  femm/materials/specimen/.style={mu r=1,conductivity=58000000}}}
\end{codeexample}
The two |\pgfmathsetmacro| calls evaluate expressions before passing their
numeric values to the mesh keys. The |ndt/| styles are local example styles,
not another installed package interface. They can be edited and composed with
normal PGF |.style| and |.append style| handlers.

|type=axisymmetric| is essential here. It solves the finite probe and its
surrounding air. The |type=axial| long-solenoid reduction of
Section~\ref{sec:currents} prescribes an axial field and has no probe terminal
impedance to measure. The outer boundary in this tutorial is the usual zero
vector-potential truncation. Enlarging the air domain and refining the wall
mesh are separate checks.

\subsection{Paint copper, remove metal with air, add the winding}
Kamil first paints a locally refined air band, then the complete copper wall.
The order matters: painting the air band last would erase the specimen in
that band. The groove is a later \emph{air region}, replacing part of the wall.
It is not an excluded mesh hole. Magnetic field must still be solved in the
space left by the removed metal.
\begin{codeexample}[code only]
\path[femm/region={material=air,mesh size=\NearMesh}]
  (0,-15) rectangle (12,15);
\filldraw[fill=blue!10,
  femm/region={material=specimen,mesh size=\Mesh}]
  (5,-25) rectangle (8,25);
% Omit this path for the sound reference:
\filldraw[fill=white,draw=red,
  femm/region={material=air,mesh size=\GrooveMesh}]
  (6.5,-1.5) rectangle (8,1.5);
\filldraw[fill=orange!25,
  femm/region={material=copper,circuit=probe,turns=60,
    winding fill=.6,mesh size=\Mesh}]
  (9,-1.5) rectangle (11,1.5);
\end{codeexample}
The specimen is a solid conductor and develops resolved eddy currents.
The named winding is stranded: it represents insulated fine wires carrying
the prescribed current, with homogenised copper resistance. It does not become
a massive copper ring just because its region uses |material=copper|.
The specified 1 A is a peak phasor, so the terminal power uses the factor $1/2$.
All sources follow $e^{+\mathrm i\omega t}$.

\subsection{Look at the field, then at the current}
The example draws magnetic flux lines in blue and normal-current symbols in
red. Both pictures use $\omega t=0$, so their arrows and symbols describe the
same instant. In the meridional section, positive $J_\varphi$ points into the
page and is drawn as a cross; a dot points out of the page. Its actual path is
a circle about the axis, not a loop lying in the $(r,z)$ picture.
\begin{codeexample}[code only]
\pic[femm/phase=0,femm/lines=17,
  femm/field style={blue!65!black},femm/arrow spacing=11mm]
  {femm field};
\pic[femm/currents={material=specimen,phase=0,spacing=.9,
  reference current=2000000,symbol radius=.9pt},
  femm/eddy style={red!70!black}] {femm eddy currents};
\end{codeexample}
The fixed current reference is in $\mathrm{A/m^2}$; it gives both panels the
same opacity scale. Magnetic contour levels are chosen separately in each
picture, so line counts are not a quantitative measure of the difference.
The view is clipped to a close-up while the complete tube and air domain
remain in the calculation. A TikZ |\clip| changes the visible drawing; region capture still passes
the complete, unclipped physical path to the solver.

A second comparison follows the current along a line inside the remaining
wall. The line lies at $r=6.25\,\mathrm{mm}$, below the floor of even the
$1.5\,\mathrm{mm}$ groove. Thus it never jumps into air at the defect:
\begin{codeexample}[code only]
\path[femm/profile={name=ndt-groove,samples=161}]
  (6.25,-10)--(6.25,10);
% Later, inside a PGFPlots axis:
\femmplot[profile=ndt-groove,component=Jeddy,abscissa=z,
  representation=amplitude,plot options={red!70!black,thick}]
\end{codeexample}
The pic is instantaneous, whereas the curve is a phasor amplitude. In the
complete example the axis converts $\mathrm{A/m^2}$ to $\mathrm{MA/m^2}$
with a normal PGFPlots |y filter|. The deeper groove concentrates current in
part of the remaining wall: the local curve may rise even while the integrated
specimen loss falls. There is no rule that every local current must decrease
when metal is removed. The plotted differences come from the solved fields.
\luafemmexamplebegin{tutorial-ndt-fields}

% manual example: tutorial-ndt-fields/1

\begin{codeexample}[setup code]
\definecolor{ndtblue}{RGB}{25,83,112}
\definecolor{ndtred}{RGB}{192,69,44}
\def\Mesh{.4}\def\Frequency{2000}\def\Position{0}
\tikzset{ndt/model/.style={femm/problem={type=axisymmetric,
    frequency=\Frequency,xmax=45,ymin=-45,ymax=45,mesh size=2},
    femm/circuit={name=probe,current=1},
    femm/materials/specimen/.style={mu r=1,conductivity=58000000}},
  ndt/tube/.style={fill=ndtblue!12,draw=ndtblue!55,
    femm/region={material=specimen,mesh size=\Mesh}},
  ndt/coil/.style={fill=orange!25,draw=orange!80!black,
    femm/region={material=copper,circuit=probe,turns=60,
      winding fill=.6,mesh size=\Mesh}}}
\pgfmathsetmacro{\NearMesh}{1.5*\Mesh}
\pgfmathsetmacro{\GrooveMesh}{\Mesh/2}
\end{codeexample}

The following picture uses the preceding definitions.

% manual example: tutorial-ndt-fields/2

\begin{codeexample}[width=0pt]
\foreach \Case/\Depth/\Title in {sound/0/Sound tube,groove/1.5/Annular groove}{
  \begin{tikzpicture}[ndt/model,scale=2,
      femm/cache={file=\jobname-\Case,mode=auto},font=\small]

    \path[femm/region={material=air,mesh size=\NearMesh}]
      (0,-15) rectangle (12,15);
    \begin{scope}
      \clip (-.5,-12) rectangle (19,12);
      \filldraw[ndt/tube] (5,-25) rectangle (8,25);
      \ifdim\Depth pt>0pt

        \filldraw[fill=white,draw=ndtred,thick,
          femm/region={material=air,mesh size=\GrooveMesh}]
          ({8-\Depth},-1.5) rectangle (8,1.5);
      \fi
      \filldraw[ndt/coil] (9,{\Position-1.5}) rectangle (11,{\Position+1.5});
      \path[femm/profile={name=ndt-\Case,samples=161}]
        (6.25,-10)--(6.25,10);
      \pic[femm/phase=0,femm/lines=17,femm/arrow spacing=11mm,
        femm/arrow length=1mm,femm/arrow width=.65mm,
        femm/field style={ndtblue,line width=.35pt}] {femm field};
      \pic[femm/currents={material=specimen,phase=0,spacing=.9,
        reference current=2000000,symbol radius=.9pt},
        femm/eddy style={ndtred}] {femm eddy currents};
      \draw[ndtred,densely dashed] (6.25,-10)--(6.25,10);
    \end{scope}
    \draw[gray,dash dot,-Stealth] (0,-13)--(0,13) node[above] {$z$};
    \draw[gray,-Stealth] (0,-13)--(18,-13) node[right] {$r$};
    \node[anchor=west,text=ndtblue,font=\bfseries] at (1,15) {\Title};
    \draw[orange!80!black,-Stealth] (14,6)
      node[above,fill=white,inner sep=1pt] {60-turn probe}
      --(11,1.5);
    \ifdim\Depth pt>0pt
      \draw[ndtred,-Stealth] (15,-6)
        node[below,fill=white,inner sep=1pt] {$d=1.5$ mm}--(7,-1);
    \else
      \draw[ndtblue,-Stealth] (15,-8)
        node[below,fill=white,inner sep=1pt] {Copper}--(7,-7);
    \fi
    \femmmeasure[name=terminal-\Case,quantity=circuit,circuit=probe]
  \end{tikzpicture}
}
\end{codeexample}

% manual example: tutorial-ndt-fields/3

\begin{codeexample}[width=0pt]
\begin{tikzpicture}
  \begin{axis}[width=9cm,height=5.8cm,grid=major,
    grid style={gray!15},axis line style={gray!60},
    xlabel={Axial coordinate $z$ (mm)},
    ylabel={$|\widehat J_\varphi|$ ($\mathrm{MA/m^2}$)},
    xmin=-10,xmax=10,ymin=0,no markers,
    legend style={draw=none,fill=white,font=\small},legend pos=north east,
    y filter/.code={\pgfmathparse{#1/1000000}},scaled y ticks=false]
    \femmplot[profile=ndt-sound,component=Jeddy,abscissa=z,
      representation=amplitude,plot options={ndtblue,thick}]
    \addlegendentry{Sound tube}
    \femmplot[profile=ndt-groove,component=Jeddy,abscissa=z,
      representation=amplitude,plot options={ndtred,thick,dashed}]
    \addlegendentry{1.5 mm groove}
  \end{axis}
\end{tikzpicture}
\end{codeexample}

\luafemmexampleend{tutorial-ndt-fields}

\subsection{What does the instrument measure?}
Kamil cannot put a current-density probe inside every triangle of the tube.
He can, however, measure the voltage needed to drive the coil. Its complex
impedance is
\[
 \widehat V=R_{\rm coil}\widehat I+\mathrm i\omega\widehat\lambda,
 \qquad Z=\frac{\widehat V}{\widehat I}=R+\mathrm iX.
\]
The flux linkage includes the reaction of the induced specimen currents.
The existing circuit measurement supplies both components:
\begin{codeexample}[code only]
\femmmeasure[name=probe-reading,quantity=circuit,circuit=probe]
% Expandable SI numbers, suitable for subsequent PGF calculations:
\femmresult{probe-reading}{impedance_real}
\femmresult{probe-reading}{impedance_imaginary}
\end{codeexample}
|impedance_real| is the total AC resistance seen at the terminals;
|resistance| alone is the winding's DC copper resistance. Kamil must retain
|impedance_real| to see specimen losses. For this one-coil, passive-conductor
model,
\[
 \tfrac12|\widehat I|^2\bigl(\mathop{\rm Re}Z-R_{\rm coil}\bigr)
 =P_{\rm specimen}.
\]
The automated tutorial test checks this identity independently of the drawing.
The coil itself is unchanged between the sound and grooved specimens, so
its copper resistance cancels in the difference.

This use of the coil's impedance is the usual eddy-current inspection
principle: a change in the specimen changes induced currents and their
magnetic reaction on the probe. The
\href{https://infotechnology.fhwa.dot.gov/eddy-current-testing/}{FHWA introduction
to eddy-current testing} describes both impedance-plane and distance-based
signal displays, as well as the influence of probe spacing and material
properties. Our example calculates a particular idealised groove response;
it does not import an instrument trace or impose a chosen signal shape.

\subsection{Move the probe and subtract a matched reference}
Let $s$ be the axial coordinate of the coil centre. Only the winding rectangle
moves:
\begin{codeexample}[code only]
\path[femm/region={material=copper,circuit=probe,turns=60,
  winding fill=.6,mesh size=\Mesh}]
  (9,{\Position-1.5}) rectangle (11,{\Position+1.5});
\femmmeasure[name=\Case-\Index,quantity=circuit,circuit=probe,
  parameter=\Position]
\end{codeexample}
Every position starts a new |tikzpicture| and declares its complete model.
Do not shift the winding after solving or add a groove after a field pic has
already triggered meshing. The scan is a sequence of steady harmonic states;
no velocity or motion-induced current term is used. Physically this describes
measurements after the electromagnetic transient has settled at each position.

For every $s$, solve the sound tube and both groove depths, then form
\[
 \Delta Z(s)=Z_{\rm groove}(s)-Z_{\rm sound}(s),\qquad
 \Delta R=\mathop{\rm Re}\Delta Z,\quad
 \Delta X=\mathop{\rm Im}\Delta Z.
\]
The sound reference is evaluated at the \emph{same position}. A finite tube
can change its coupling as the coil approaches an end; subtracting a single
reference value for the whole scan would mix that change into the defect
signal. This is a numerical reference subtraction, not a model of a physical
two-coil differential probe or a bridge circuit.

The full source uses a named TeX box to execute each model without printing
57 intermediate pictures. Named measurements remain available after the box
is discarded. Each case/position has its own native cache. On a later
compilation, changing graph styles or labels restores the solved phasors.
Changing physical geometry, frequency or mesh settings invalidates the
appropriate caches automatically.

The difference is small compared with the total impedance. Use PGF's
floating-point unit when calculating it, and convert to milliohms afterward:
\begin{codeexample}[code only]
\begingroup
  \pgfkeys{/pgf/fpu=true,/pgf/fpu/output format=fixed}
  \pgfmathparse{1000*(\femmresult{deep-10}{impedance_real}
    -\femmresult{sound-10}{impedance_real})}
  \xdef\CentreDeltaR{\pgfmathresult}
\endgroup
\end{codeexample}
In the distributed position list, index 10 is $s=0$. The floating-point keys
are deliberately grouped: they must not change the arithmetic used by later
TikZ geometry. The complete example collects coordinate lists with ordinary
TeX macros and draws them with |\addplot|. No hand-written Lua is necessary.
\luafemmexamplebegin{tutorial-ndt-scan}

Kamil first names the mesh settings, the specimen and the two output channels.
\begin{codeexample}[setup code]
\definecolor{ndtblue}{RGB}{25,83,112}
\definecolor{ndtred}{RGB}{192,69,44}
\def\Mesh{.4}\def\Frequency{2000}
\newbox\NDTCalculation
\def\ShallowR{}\def\ShallowX{}\def\ShallowPlane{}
\def\DeepR{}\def\DeepX{}\def\DeepPlane{}
\tikzset{ndt/model/.style={femm/problem={type=axisymmetric,
    frequency=\Frequency,xmax=45,ymin=-45,ymax=45,mesh size=2},
    femm/circuit={name=probe,current=1},
    femm/materials/specimen/.style={mu r=1,conductivity=58000000}}}
\end{codeexample}
The following styles affect the axes only. The local mesh lengths belong to the physical model.
\begin{codeexample}[setup code]
\pgfplotsset{ndt/axis/.style={grid=major,grid style={gray!15},
    axis line style={gray!60},tick label style={font=\small},
    legend style={draw=none,fill=white,font=\small},
    scaled ticks=false},
  ndt/shallow/.style={ndtblue,thick,mark=*,mark size=1.3pt},
  ndt/deep/.style={ndtred,thick,mark=square*,mark size=1.5pt}}
\pgfmathsetmacro{\NearMesh}{1.5*\Mesh}
\pgfmathsetmacro{\GrooveMesh}{\Mesh/2}
\end{codeexample}
A command describes one probe position. Its arguments come from the scan loop below; the discarded box suppresses intermediate drawings, not the calculation.
\begin{codeexample}[setup code]
\newcommand\ScanProbe{%
\setbox\NDTCalculation=\hbox{
      \tikzpicture[ndt/model,
          femm/cache={file=\jobname-\Case-\Index,mode=auto}]
        \path[femm/region={material=air,mesh size=\NearMesh}]
          (0,-15) rectangle (12,15);
        \path[femm/region={material=specimen,mesh size=\Mesh}]
          (5,-25) rectangle (8,25);
        \ifdim\Depth pt>0pt
          \path[femm/region={material=air,mesh size=\GrooveMesh}]
            ({8-\Depth},-1.5) rectangle (8,1.5);
        \fi
        \path[femm/region={material=copper,circuit=probe,turns=60,
            winding fill=.6,mesh size=\Mesh}]
          (9,{\Position-1.5}) rectangle (11,{\Position+1.5});
        \femmmeasure[name=\Case-\Index,quantity=circuit,circuit=probe,
          parameter=\Position]
      \endtikzpicture}}
\end{codeexample}
A second command subtracts the sound-tube reference at the same position. The floating-point settings are local to this subtraction.
\begin{codeexample}[setup code]
\newcommand\RecordSignals{%
\begingroup\pgfkeys{/pgf/fpu=true,/pgf/fpu/output format=fixed}
    \foreach \Case/\Prefix in {shallow/Shallow,deep/Deep}{
      \pgfmathparse{1000*(\femmresult{\Case-\Index}{impedance_real}
        -\femmresult{sound-\Index}{impedance_real})}
      \edef\RealSignal{\pgfmathresult}
      \expandafter\xdef\csname\Prefix R\endcsname{
        \csname\Prefix R\endcsname (\Position,\RealSignal)}
      \pgfmathparse{1000*(\femmresult{\Case-\Index}{impedance_imaginary}
        -\femmresult{sound-\Index}{impedance_imaginary})}
      \expandafter\xdef\csname\Prefix X\endcsname{
        \csname\Prefix X\endcsname (\Position,\pgfmathresult)}
      \expandafter\xdef\csname\Prefix Plane\endcsname{
        \csname\Prefix Plane\endcsname (\RealSignal,\pgfmathresult)}
    }
  \endgroup}
\end{codeexample}
The loop now reads like the experiment: place the coil, compute each specimen, and retain the two differences. Each point is an independent harmonic solve.
\begin{codeexample}[setup code]
\foreach \Position [count=\Index] in
  {-12,-10,-8,-6,-5,-4,-3,-2,-1,0,1,2,3,4,5,6,8,10,12}{
  \foreach \Case/\Depth in {sound/0,shallow/.5,deep/1.5}{
    \ScanProbe
  }
  \RecordSignals
}
\end{codeexample}
The resistive channel compares the change of input loss at fixed current.
\begin{codeexample}[width=0pt]
\begin{tikzpicture}
  \begin{axis}[ndt/axis,width=8cm,height=6cm,xmin=-12,xmax=12,
      xlabel={Probe centre $s$ (mm)},ylabel={$\Delta R$ (m$\Omega$)},
      title={Resistive channel},legend pos=south east]
    \addplot[gray,densely dashed,forget plot] coordinates {(-12,0)(12,0)};
    \addplot[ndt/shallow] coordinates {\ShallowR};
    \addlegendentry{$d=0.5$ mm}
    \addplot[ndt/deep] coordinates {\DeepR};
    \addlegendentry{$d=1.5$ mm}
  \end{axis}
\end{tikzpicture}
\end{codeexample}
The reactive channel uses the imaginary part of the same impedance difference.
\begin{codeexample}[width=0pt]
\begin{tikzpicture}
  \begin{axis}[ndt/axis,width=8cm,height=6cm,xmin=-12,xmax=12,
      xlabel={Probe centre $s$ (mm)},ylabel={$\Delta X$ (m$\Omega$)},
      title={Reactive channel}]
    \addplot[gray,densely dashed,forget plot] coordinates {(-12,0)(12,0)};
    \addplot[ndt/shallow] coordinates {\ShallowX};
    \addplot[ndt/deep] coordinates {\DeepX};
  \end{axis}
\end{tikzpicture}
\end{codeexample}
Finally, plot the two channels against one another. These points describe probe positions, not temporal phases.
\begin{codeexample}[width=0pt]
\begin{tikzpicture}
  \begin{axis}[ndt/axis,width=10cm,height=7cm,
      xlabel={$\Delta R$ (m$\Omega$)},ylabel={$\Delta X$ (m$\Omega$)},
      title={The same scan in the impedance plane},legend pos=north east]
    \addplot[ndt/shallow] coordinates {\ShallowPlane};
    \addlegendentry{$d=0.5$ mm}
    \addplot[ndt/deep] coordinates {\DeepPlane};
    \addlegendentry{$d=1.5$ mm}
    \node[anchor=north east,font=\small] at (axis cs:0,0) {Far from groove};
  \end{axis}
\end{tikzpicture}
\end{codeexample}
\luafemmexampleend{tutorial-ndt-scan}

\subsection{Read the three graphs together}
The first graph is the resistive channel. For these parameters it has a
negative central response and smaller positive shoulders. The middle of the
groove gives lower total specimen loss, while the off-centre response also
contains redistribution and changed coupling. A generic ``defect signal''
should not be substituted for this computed shape.

The reactive channel has a positive central response. The deeper groove has
a larger complex signal, although doubling groove depth would not in general
double either component. The impedance-plane graph uses the same computed
points, with $\Delta R$ horizontally and $\Delta X$ vertically. It does not
plot the temporal sine wave. The geometry is symmetric under $s\mapsto-s$,
so the two halves of the scan nearly retrace the same trajectory. Small
separations between them reflect the discretisation, not hysteresis.

The markers identify actual solves. Lines join adjacent samples without
claiming an interpolated extremum or adding a fitted defect model. Use more
positions near a narrow response to improve spatial sampling; use more
triangles to improve field accuracy. These are different operations.

\subsection{Before interpreting a small difference}
Kamil has three independent resolution controls:
\begin{enumerate}
\item Reduce local |mesh size| in the metal, groove and nearby air. Compare
$\Delta Z$ itself: a small change in the much larger total $Z$ does not by
itself establish convergence of the difference.
\item Enlarge the outer air domain while retaining the same local mesh
settings. The zero-potential boundary is a truncation, not a real shield.
For this example, the initially tried $25\,\mathrm{mm}$ radial extent changed
the central differential signal by about $2.7\%$ when enlarged to $45$ mm,
which is why the complete sources use the larger domain.
\item Refine the list of probe positions independently of the field mesh.
The source places extra points near the groove, where the signal changes most.
\end{enumerate}
The regression checks the two depth responses, the power
identity, mesh refinement, outer-domain enlargement, scan symmetry and the
small far-from-groove signal. Its recorded results and tolerances are in
\nolinkurl{doc/validation.md}. They establish numerical sensitivity for this
teaching geometry, not an experimental detection threshold.

Finally, a change in impedance is not uniquely a change in groove depth.
Conductivity, frequency and coil clearance also change the coupling. In a
real inspection these require calibration; the present calculation includes
no instrument noise or uncertainty distribution. In an encircling-coil model,
increasing the coil radius changes its circumference and therefore its copper
resistance as well as its clearance. Keep that distinction when extending the
exercise. Lateral wobble and isolated cracks require a three-dimensional model.
The useful result here is a reproducible chain from a declared material defect
to fields, currents and a measurable complex coil response.

\clearpage
